\section{Validation, scope, and a practical continuation order}

\subsection{What was checked in this package}

The standalone Lean files \texttt{FiniteCore.lean},
\texttt{Targets.lean}, and \texttt{FiniteBranch.lean} were compiled
with Lean 4.32.0. The first
contains actual proofs of the bounded search, energy, endpoint, and
contract lemmas. The second fixes concrete definitions and propositions
for future proofs. Its successful compilation establishes that those
definitions and statements are well-typed, not that the target
propositions have been proved. The archive retains the compiler output
and its scope declaration. The third supplies Gaussian-rational
algorithms and five kernel-checked coefficient computations; its
general analytic goal fields remain uninstantiated.

The standard-library Python reference program uses exact
\texttt{Fraction} arithmetic and checks the following.
\begin{itemize}
\item All $256$ binary sequences of length eight, extended constantly,
are fed to the finite-variation selector with their actual variation
budget. Nonmonotone examples are included. Two inputs violating the
finite budget hypotheses are rejected. A two-sequence example with an
exceptional set verifies the synchronized output and the bound $2046$.
\item The simple branch of $Y^2=1+X$ is constructed through order $20$
by truncated polynomial iteration. Its coefficients agree with the
independent binomial recurrence. The run checks $21$ weighted bounds,
$231$ stabilization comparisons, $924$ finite-tail inequalities, and
$84$ residual inequalities, all exactly.
\item The sharp four-point hyperelliptic example and a second example
with $B\ne0$ verify the curve equations, norm polynomial, repeated
synthetic divisions, and point count.
\item All eight border identities for the size-four quotient certificate
are verified by coefficient comparison. The evaluation matrix on the
four displayed solutions has determinant $-4$, confirming that the
example attains the dimension bound.
\end{itemize}
These are implementation and example checks. The proofs of the universal
mathematical assertions are in the manuscript; their full machine
formalization remains as described above. In particular, no numerical
test has been used as evidence for an unresolved global height estimate.

The mathematical chapters were also reviewed independently during
preparation. The review checked the degree and multiplicity argument,
the finite algebra dimension contradiction, the common-window endpoint
conventions, the analytic norm and precision constants, and the
quantifier-free induction boundaries. The report was compiled to PDF
and its rendered pages inspected. The archive records reproducibility
commands, source provenance, and file checksums.

\subsection{Unresolved mathematical dependencies}

The proposed framework does not establish the following stronger claims:
\begin{enumerate}
\item a uniform primitive recursive bound for all rational points of
arbitrary genus-at-least-two curves, or for integral points on arbitrary
curves covered by Siegel's theorem;
\item a uniform generator algorithm for arbitrary Mordell--Weil groups
from abstract finiteness or Selmer information alone;
\item the finite-tuple quantitative Subspace Theorem required by the
conditional hyperelliptic point-count corollary;
\item a general effective Roth cutoff, or all of the analytic estimates
and uniform witnesses left open by the prime-distribution packet;
\item uniform exact singular Puiseux coefficient selection from arbitrary
Cauchy presentations, or a generic zero test for such presentations;
\item a global Cauchy modulus extracted from a finite-window selector;
\item an implemented automatic compiler from arbitrary Lean proofs to iPRA.
\end{enumerate}
The first four entries are left open here. Their absence is not a proof
of unprovability. The fifth and sixth contain precise representational
or computability obstructions, explained in the inspected source and
in this report. The last entry would require much more than the
restricted proof-construction discipline established here.

\subsection{Recommended implementation order}

The best first complete embedded theorem is the hyperelliptic fiber
bound. Its data are exact rationals, its new idea is isolated, and its
induction invariants are short. It also produces an immediate usable
contract for an existing research route. The size-four example supplies
a small end-to-end test of both the norm and the finite algebra paths.

Next implement polynomial certificate soundness and the spanning
certificate theorem. These modules can serve a broad range of algebraic
and geometric packets. In parallel, the bounded selector and
common-window theorem can be translated into the object theory using
only finite sums and counting.

The analytic branch theorem should then reuse the same finite polynomial
arithmetic, adding weighted norms and precision contracts. Its key
advantage is that the perturbation bound is independent of iteration
count. This provides a stable interface for later analytic applications,
instead of requiring each application to rediscover how precision grows
through a long iteration.

The long-term integration point is a collection of reusable, proved
finite contracts with explicit demand functions. Each existing packet
can expose which of these contracts it supplies, which it consumes,
and which it still lacks. That organization makes further mathematical
progress concrete: a missing quantitative theorem remains visible,
while previously established finite algebra and analytic estimates
become reusable derivations.
